\documentclass{article}
\usepackage[utf8]{inputenc}
\usepackage[T1]{fontenc}
\usepackage{geometry}
\usepackage{graphicx}
\usepackage{float}
\usepackage{hyperref}
\usepackage[english,indonesian]{babel}

\title{Optimasi Pemrosesan Template LaTeX Menggunakan Pendekatan Berbasis AI}
\author{Academic Support Agent}
\date{\today}

\begin{document}
\maketitle

\begin{abstract}

\end{abstract}

\chapter{PENDAHULUAN}


\chapter{TINJAUAN PUSTAKA}


% Original content (commented out):
% \documentclass{article}
\usepackage[utf8]{inputenc}

\title{Resumo: Capítulo 7}
\author{aramigalene }
\date{July 2019}

\begin{document}

\maketitle

\section{Introduction}


O WSMO (Web Service Modeling Ontology) é uma ontologia formal que modela os serviços da web, que utiliza  A linguagem WSML ( Web Service Modeling Language) e um ambiente de execução WSMX ( Web Service Modeling Enviroment). Do WSMF foram adotados quatro principasi elementos para descrever os serviços da web: Ontologias, Goals (definem os problemas a serem resolvidos), descrições (dos serviços) e os mediadores (possibilitam a interoperabilidade na resolução de problemas).
\section{Tecnologias}
As tecnologias de serviço web atuais exigem ainda uma interação manual do programador para encontrar serviços específicos da web wue atendam as suas necessidades, o que limita a escalabilidade e aumenta os custos. Para automatizar tarefas relacionadas aos serviços web aplica-se o WSMO.
\section{Princípios de Design}
O WSMO especifica os seguintes princípios de design: É baseado em ontologias (descrevendo todos os dados), desacoplamento rigoroso (recursos definidos de forma isolada e independente), conformidade  com a web (Inerente ao URI), centralidade da mediação (Muito importante para a boa implantação dos serviços web), separação de ontologias (De acordo com os contextos específicos), distingue descrição (Dos serviços) de execução (Tecnologias executáveis), semântica de execução (define bem a funcionalidade e o comportamento dos sistemas) e difere serviços da web (entidade computacional) de serviço (valor real).
\section{MOF}
O WSMO utiliza o MOF (Meta Object Facility), recurso que define seu meta modelo, fornece uma estrutura que especifica meta modelos neutros e descreve semanticamente os serviços web de forma separada. Um dos padrões mais aplicados é o conjunto de metadados Dublin Core (DC) para recursos de informações entre dois domínios.
\section{Mediadores}
São utilizados para possibilitar a união entre ontologias nomadamente incompatíveis, realizando uma tradução de uma para outra.
\section{Conceito}
Representam a terminologia básica de um domínio composto por atributos como nomes e tipos distintos. Um conceito pode ser um sub-conceito (parte de um super-conceito) ou um super-conceito (formado de sub-conceitos). Em WSMO a intrução hasFunction é uma relação especial que possui um intervalo unário e especifica o domínio junto com os parâmetros. A expressão lógica hasDefinition possibilita capturar a semântica de uma relação ou função.
\section{Instância}
Em WSMO é possível instanciar conceitos e relações, atribuindo-se valores aos atributos e parâmetros. Instâncias podem ser explícitas a ontologia, porém, se a ontologia tiver muitas instâncias geralmente é utilizado um link que possibilite o armazenamento externo.
\section{Axiomas}
Um axioma é definido como uma expressão lógica com propriedades não funcionais que fornece mecanismos para refinar conceitos, relações ou definições das funções, expressar restrições e adicionar diferentes conhecimentos de domínio axiomático.
\section{Web Services}
São entidades computacionais que fornecem alguma funcionalidade e associam valores a um determinado domínio. Em WSMO é uma funcionalidade disponível via método importado da notação MOF. As ontologias que definem os serviços podem ser importados por meio dos métodos importsOntology e usesMediator.
\section{}

\end{document}


\end{document}